%% Drafted from the mapped corpus by scripts/draft_topics.py.
% Model: gpt-5.6-terra; source characters: 17554.
\section{Множества и аксиоматичен подход}

Понятието \emph{множество} се приема за първично: използва се без дефиниция. Елементите на множество $A$ се означават с $x\in A$, а твърдението, че $x$ не е негов елемент, с $x\notin A$. \begin{defn}[Подмножество]
Множеството $A$ е подмножество на $B$, означено с $A\subseteq B$, когато всеки елемент на $A$ е елемент и на $B$:
\[
A\subseteq B\quad\Longleftrightarrow\quad
(\forall x)\,(x\in A\Rightarrow x\in B).
\]
\end{defn}

\begin{keyaxiom}[За обема]
Две множества $A$ и $B$ са равни, ако и само ако имат едни и същи елементи:
\[
A=B\quad\Longleftrightarrow\quad(\forall x)\,(x\in A\Leftrightarrow x\in B).
\]
Формалният запис е
\[
(\forall X)(\forall Y)\left(
(\forall z)\,(z\in X\Leftrightarrow z\in Y)\Rightarrow X=Y
\right).
\]
\end{keyaxiom}

\subsection{Аксиоми за построяване на множества}

\begin{keyscheme}[За отделянето]
Нека $A$ е множество и $\varphi(x,u_1,\ldots,u_n)$ е свойство. Тогава съществува множество, състоящо се точно от онези елементи на $A$, които удовлетворяват това свойство:
\[
B=\{x\in A\mid \varphi(x,u_1,\ldots,u_n)\}.
\]
Формално:
\[
(\forall u_1)\cdots(\forall u_n)(\forall A)(\exists B)(\forall x)
\left(x\in B\Leftrightarrow
\bigl(x\in A\land\varphi(x,u_1,\ldots,u_n)\bigr)\right).
\]
\end{keyscheme}

Схемата за отделянето позволява да се образуват подмножества на вече дадено множество чрез зададено условие. Например множеството на четните естествени числа, които са не по-големи от $20$, може да се запише като
\[
\{x\in\mathbb N\mid x\leq 20\text{ и }x\text{ е четно}\}.
\]

\begin{keyaxiom}[За степенното множество]
За всяко множество $A$ съществува множество от всички негови подмножества, наречено \emph{степенно множество} на $A$ и означавано с $\mathcal P(A)$:
\[
\mathcal P(A)=\{X\mid X\subseteq A\}.
\]
\end{keyaxiom}

Например
\[
\mathcal P(\varnothing)=\{\varnothing\},
\qquad
\mathcal P(\{a,b\})=
\{\varnothing,\{a\},\{b\},\{a,b\}\}.
\]
За крайно множество $A$ с $|A|=n$ е валидно
\[
|\mathcal P(A)|=2^{|A|}=2^n.
\]

\begin{keyaxiom}[За индуктивно генерираните множества]
Нека $U$ е множество, $\varnothing\ne M_0\subseteq U$, а $F$ е множество от операции $f:U^{k_f}\to U$ с крайна положителна арност $k_f$. Тогава множеството $M$ се строи по следния начин:

\begin{enumerate}
\item \emph{База:} $M_0\subseteq M$, т.е. в $M$ са всички елементи на $M_0$.
\item \emph{Индуктивно предположение:} нека $M_i$ е вече построеното на $i$-тата стъпка множество. То е непразно, защото $\varnothing\neq M_0\subseteq M_i$.
\item \emph{Индуктивна стъпка:} в $M_{i+1}$ включваме всички досегашни елементи и всички елементи, които се получават от тях чрез прилагане на операциите от $F$.
\item \emph{Заключение:} в $M$ няма други елементи освен базовите и получените чрез краен брой индуктивни стъпки.
\end{enumerate}

При едноместни операции конструкцията се записва формално като
\[
M_{i+1}=M_i\cup\{f(m)\mid f\in F,\ m\in M_i\},
\qquad
M=\bigcup_{i=0}^{\infty}M_i.
\]
При многоместна операция $f$ се прилагат всички допустими наредени групи от елементи на $M_i$. Построеното множество се означава с
\[
M:=\langle M_0,F\rangle.
\]
\end{keyaxiom}

Заключението изразява минималността на построеното множество: елементите му са точно базовите елементи и елементите, достижими чрез краен брой приложения на допустимите операции.

\begin{supp}[Бележка]
Обединението с $M_i$ в индуктивната стъпка е съществено: то запазва вече получените елементи при преминаването към следващото ниво.
\end{supp}

\section{Математическа индукция}

Математическата индукция е метод за доказване на твърдения за естествени числа. Нека $m\in\mathbb N$ и $\varphi(n)$ е твърдение, зависещо от $n$.

\begin{keythm}[Принцип на слабата математическа индукция]
Ако са изпълнени:

\begin{enumerate}
\item базата: $\varphi(m)$ е вярно;
\item индуктивната стъпка:
\[
(\forall k\geq m)\,\bigl(\varphi(k)\Rightarrow\varphi(k+1)\bigr),
\]
\end{enumerate}

то
\[
(\forall n\in\mathbb N,\ n\geq m)\,\varphi(n).
\]
\end{keythm}

\begin{proof}
Използваме принципа за най-малкия елемент на естествените числа. Ако има $n\ge m$, за което $\varphi(n)$ е невярно, нека $n_0$ е най-малкото такова число. Базата дава $n_0>m$. От минималността $\varphi(n_0-1)$ е вярно, а индуктивната стъпка тогава дава $\varphi(n_0)$ --- противоречие. Следователно няма контрапример. Принципът за най-малкия елемент и индукцията са еквивалентни основни свойства на $\mathbb N$; тук първият е приет за изходен.
\end{proof}

Прилагането на принципа има ясен строеж. Първо се доказва твърдението за началната стойност. След това се приема, че е вярно за произволно, но фиксирано $k$, и при това допускане се доказва за $k+1$. Така от верността за $m$ следва верността за $m+1$, после за $m+2$ и т.н.

\begin{keythm}[Принцип на пълната математическа индукция]
Ако $\varphi(m)$ е вярно и за всяко $k>m$ от
\[
\varphi(m)\land\varphi(m+1)\land\cdots\land\varphi(k-1)
\]
следва $\varphi(k)$, то $\varphi(n)$ е вярно за всяко естествено число $n\geq m$.
\end{keythm}

\begin{proof}
Прилагаме слабата индукция към $\psi(k)\equiv\bigwedge_{j=m}^{k}\varphi(j)$. Базата $\psi(m)$ е дадена. Ако $\psi(k)$ е вярно, предпоставката на пълната индукция дава $\varphi(k+1)$, следователно и $\psi(k+1)$. Получаваме $\psi(k)$ за всяко $k\ge m$, откъдето следва $\varphi(k)$.
\end{proof}

Разликата е в индуктивното предположение: при слабата индукция се използва само $\varphi(k)$, а при пълната индукция може да се използва верността на твърдението за всички предишни стойности.

\begin{example}
Да се докаже, че за всяко $n\in\mathbb N$:
\[
0+1+\cdots+n=\frac{n(n+1)}{2}.
\]

За $n=0$ равенството е вярно. Нека е вярно за $n=k$. Тогава
\[
0+1+\cdots+k+(k+1)
=
\frac{k(k+1)}{2}+(k+1)
=
\frac{(k+1)(k+2)}{2}.
\]
Следователно формулата е вярна и за $k+1$, а по математическа индукция --- за всяко $n\in\mathbb N$.
\end{example}

\section{Операции върху множества}

\begin{defn}[Операции върху множества]
Нека $A$ и $B$ са множества. Основните операции са следните:
\[
A\cup B=\{x\mid x\in A\lor x\in B\},
\]
\[
A\cap B=\{x\mid x\in A\land x\in B\},
\]
\[
A\setminus B=\{x\mid x\in A\land x\notin B\},
\]
\[
A\triangle B=(A\setminus B)\cup(B\setminus A).
\]
\end{defn}

Обединението съдържа елементите, които принадлежат поне на едно от множествата; сечението съдържа общите им елементи; разликата $A\setminus B$ съдържа елементите на $A$, които не са в $B$. Симетричната разлика съдържа елементите, които принадлежат точно на едно от двете множества.

\begin{defn}[Допълнение]
При фиксиран универсум $U$, съдържащ разглежданите множества, \emph{допълнение} на $A$ до $U$ е
\[
\overline A=U\setminus A=\{x\in U\mid x\notin A\}.
\]
\end{defn}

\begin{defn}[Обобщено обединение и сечение]
Нека $A_i\subseteq U$ за всяко $i\in I$, където $U$ е фиксиран универсум. Използват се обобщенията
\[
\bigcup_{i\in I}A_i
=
\{x\mid (\exists i\in I)\ x\in A_i\},
\]
\[
\bigcap_{i\in I}A_i
=
\{x\in U\mid (\forall i\in I)\ x\in A_i\}.
\]
При $I=\varnothing$ обединението е $\varnothing$, а сечението е $U$.
\end{defn}

\begin{prop}[Основни свойства]
За произволни множества $A$, $B$ и $C$ са валидни:

\begin{align*}
A\cup A&=A, & A\cap A&=A;\\
A\cup B&=B\cup A, & A\cap B&=B\cap A;\\
(A\cup B)\cup C&=A\cup(B\cup C), &
(A\cap B)\cap C&=A\cap(B\cap C);\\
A\cup(B\cap C)&=(A\cup B)\cap(A\cup C),&
A\cap(B\cup C)&=(A\cap B)\cup(A\cap C);\\
A\cup\varnothing&=A, & A\cap U&=A;\\
A\cup U&=U, & A\cap\varnothing&=\varnothing;\\
A\cup\overline A&=U, & A\cap\overline A&=\varnothing;\\
\overline{\overline A}&=A. &&
\end{align*}

Освен това са в сила законите на Де Морган:
\[
\overline{A\cup B}=\overline A\cap\overline B,
\qquad
\overline{A\cap B}=\overline A\cup\overline B.
\]
\end{prop}

\begin{proof}
Фиксираме произволен $x\in U$ и означаваме $p\equiv x\in A$, $q\equiv x\in B$, $r\equiv x\in C$. Принадлежността на обединение, сечение и допълнение се превежда съответно като $\lor$, $\land$ и $\neg$. Идемпотентността, комутативността, асоциативността и дистрибутивността следват от едноименните булеви тъждества (проверими за всички осем оценки на $p,q,r$). Принадлежността на $\varnothing$ е неистина, а на $U$ --- истина, което дава четирите равенства с тези множества. Законите за допълнението следват от $p\lor\neg p=1$, $p\land\neg p=0$ и $\neg\neg p=p$. Накрая $\neg(p\lor q)\equiv\neg p\land\neg q$ и $\neg(p\land q)\equiv\neg p\lor\neg q$ дават двата закона на Де Морган. Във всеки случай двете страни имат едни и същи елементи и аксиомата за обема дава равенството.
\end{proof}

\begin{example}
Нека
\[
A=\{x\in\mathbb N\mid x>1\},
\qquad
B=\{x\in\mathbb N\mid x>3\}.
\]
Тогава $B\subseteq A$ и следователно
\[
A\cup B=A,\qquad A\cap B=B,
\]
\[
A\setminus B=\{x\in\mathbb N\mid 1<x\leq 3\},
\qquad
B\setminus A=\varnothing,
\]
\[
A\triangle B=\{x\in\mathbb N\mid 1<x\leq 3\}.
\]
\end{example}

\section{Наредени двойки и декартови произведения}

\begin{keydefn}[Наредена двойка]
Наредена двойка от елементите $a$ и $b$ се означава с $(a,b)$. Тя се определя чрез характеристичното свойство
\[
(a,b)=(c,d)\quad\Longleftrightarrow\quad a=c\land b=d.
\]
\end{keydefn}

Редът е съществен: по принцип $(a,b)\neq(b,a)$. Една теоретико-множествена реализация на наредена двойка е конструкцията на Куратовски:
\[
(a,b)=\bigl\{\{a\},\{a,b\}\bigr\}.
\]

\begin{keydefn}[Наредена $n$-орка]
Полагаме $(a_1)=a_1$ като база на рекурсивния запис. За $n\geq 2$ наредената $n$-орка от $a_1,\ldots,a_n$ се означава с
\[
(a_1,\ldots,a_n)
\]
и може рекурсивно да се разбира като
\[
(a_1,(a_2,\ldots,a_n)).
\]
\end{keydefn}

\begin{keydefn}[Декартово произведение]
За множества $A$ и $B$ тяхното декартово произведение е
\[
A\times B=\{(a,b)\mid a\in A\land b\in B\}.
\]
\end{keydefn}

По принцип $A\times B\neq B\times A$, защото се променя редът на компонентите. Например
\[
\{1,2\}\times\{a\}=\{(1,a),(2,a)\},
\]
докато
\[
\{a\}\times\{1,2\}=\{(a,1),(a,2)\}.
\]

\begin{keydefn}[Обобщено декартово произведение]
За множества $A_1,\ldots,A_n$:
\[
\prod_{i=1}^{n}A_i
=
A_1\times\cdots\times A_n
=
\{(a_1,\ldots,a_n)\mid a_i\in A_i,\ 1\leq i\leq n\}.
\]
\end{keydefn}

\section{Релации}

\begin{keydefn}[$n$-местна релация]
Нека $A_1,\ldots,A_n$ са множества. Всяко подмножество
\[
R\subseteq A_1\times A_2\times\cdots\times A_n
\]
се нарича $n$-местна, или $n$-арна, релация над домейните $A_1,\ldots,A_n$.
\end{keydefn}

При $n=2$ говорим за бинарна релация. Ако домейните са еднакви, например $R\subseteq A\times A$, релацията е хомогенна. Вместо $(a,b)\in R$ често се пише $aRb$.

\begin{defn}[Домейн и множество от стойности на релация]
За $R\subseteq A\times B$ се дефинират:
\[
\operatorname{dom}(R)=
\{a\in A\mid(\exists b\in B)\ (a,b)\in R\},
\]
\[
\operatorname{rng}(R)=
\{b\in B\mid(\exists a\in A)\ (a,b)\in R\}.
\]
Това са съответно домейнът и множеството от стойности на релацията.
\end{defn}

\subsection{Свойства на бинарните релации}

\begin{defn}[Свойства на бинарна релация]
Нека $R\subseteq A\times A$. Казваме, че $R$ е:

\begin{itemize}
\item \emph{рефлексивна}, ако
\[
(\forall x\in A)\ xRx;
\]

\item \emph{антирефлексивна}, или \emph{ирефлексивна}, ако
\[
(\forall x\in A)\ \neg xRx;
\]

\item \emph{симетрична}, ако
\[
(\forall x,y\in A)\ (xRy\Rightarrow yRx);
\]

\item \emph{антисиметрична}, ако
\[
(\forall x,y\in A)\ (xRy\land yRx\Rightarrow x=y);
\]

\item \emph{силно антисиметрична}, ако за всеки две различни $x,y\in A$ е изпълнено точно едно от $xRy$ и $yRx$;

\item \emph{асиметрична}, ако
\[
(\forall x,y\in A)\ (xRy\Rightarrow\neg yRx);
\]

\item \emph{транзитивна}, ако
\[
(\forall x,y,z\in A)\ (xRy\land yRz\Rightarrow xRz).
\]
\end{itemize}
\end{defn}

Антисиметричността не означава, че за различни елементи винаги съществува сравнение. Тя забранява едновременното наличие на $xRy$ и $yRx$ при $x\neq y$. Силната антисиметричност добавя изискването за сравнимост на всяка двойка различни елементи.

\subsection{Еквивалентност и класове на еквивалентност}

\begin{keydefn}[Релация на еквивалентност]
Релация $R\subseteq A\times A$ е релация на еквивалентност, ако е рефлексивна, симетрична и транзитивна.
\end{keydefn}

\begin{keydefn}[Клас на еквивалентност]
Нека $R$ е релация на еквивалентност над $A$ и $a\in A$. Класът на еквивалентност на $a$ е
\[
[a]_R=\{b\in A\mid aRb\}.
\]
Когато релацията е ясна от контекста, се пише просто $[a]$.
\end{keydefn}

\begin{thm}
Всяка релация на еквивалентност над множество $A$ поражда разбиване на $A$ на класове на еквивалентност.
\end{thm}

\begin{proof}
По рефлексивност $a\in[a]_R$, затова класовете са непразни и покриват $A$. Ако $c\in[a]_R\cap[b]_R$, то $aRc$ и $bRc$, откъдето по симетричност и транзитивност $aRb$. За всяко $x\in[a]_R$ имаме $bRaRx$, следователно $x\in[b]_R$. Обратното включване е аналогично. Значи два класа или съвпадат, или са непресичащи се, както изисква определението за разбиване.
\end{proof}

\begin{defn}[Разбиване на множество]
Фамилията $\{A_i\mid i\in I\}$ е \emph{разбиване} на $A$, ако всички $A_i$ са непразни, обединението им е $A$ и различните части нямат общи елементи:
\[
\bigcup_{i\in I}A_i=A,
\]
\[
(\forall i,j\in I)\,
(A_i\cap A_j\neq\varnothing\Rightarrow A_i=A_j).
\]
\end{defn}

\begin{example}
В $\mathbb Z$ задаваме $aRb$, ако $a$ и $b$ дават еднакъв остатък при деление на $3$. Това е релация на еквивалентност. Класовете са числата с остатък $0$, $1$ и $2$ при деление на $3$.
\end{example}

\subsection{Частични и пълни наредби}

\begin{keydefn}[Частична наредба]
Релация $R\subseteq A\times A$ е нестрога частична наредба, ако е рефлексивна, антисиметрична и транзитивна.
\end{keydefn}

\begin{defn}[Строга частична наредба]
Строга частична наредба върху $A$ е антирефлексивна и транзитивна релация $R\subseteq A\times A$. Тези условия влекат асиметричност: от $xRy$ и $yRx$ по транзитивност би следвало $xRx$, противоречие. Допълнителната антисиметричност в изходната формулировка е излишна.
\end{defn}

\begin{keydefn}[Пълна наредба]
Частична наредба $R$ над $A$ е пълна, ако всеки два различни елемента са сравними:
\[
(\forall a,b\in A)\,
\left(a\neq b\Rightarrow(aRb\lor bRa)\right).
\]
\end{keydefn}

Релациите $\leq$ и $<$ над естествените числа са примери за пълни наредби в нестрогия и строгия вариант.

\begin{keydefn}[Минимален и максимален елемент]
Нека $R$ е частична наредба над $A$. Елементът $a\in A$ е минимален, ако
\[
(\forall b\in A)\,(b\neq a\Rightarrow\neg bRa).
\]
Той е максимален, ако
\[
(\forall b\in A)\,(b\neq a\Rightarrow\neg aRb).
\]
\end{keydefn}

\begin{thm}
Всяка частична наредба над непразно крайно множество притежава минимален и максимален елемент.
\end{thm}

\begin{proof}
Избираме $a_0\in A$. Ако той не е минимален, има различен $a_1$ с $a_1Ra_0$. Продължаваме, докато достигнем минимален елемент. Повторение $a_i=a_j$ с $i<j$ е невъзможно: транзитивността и антисиметричността биха дали $a_i=a_{i+1}$, въпреки строгия избор. Понеже $A$ е крайно, процесът спира. Прилагаме същото разсъждение към обратната наредба $R^{-1}$ и получаваме максимален елемент.
\end{proof}

\begin{keydefn}[Диаграма на Хассе]
За крайно частично наредено множество $(A,\le)$ диаграмата на Хассе има върхове елементите на $A$. Ребро от $x$ нагоре към $y$ се чертае точно когато $x<y$ и няма $z$ с $x<z<y$. Тогава $y$ покрива $x$. Рефлексивните примки и ребрата, следващи по транзитивност, не се рисуват. Това уточнява общото описание на граф на релация в източника.
\end{keydefn}

При частична наредба тя позволява удобно да се видят несравнимите елементи, както и минималните и максималните елементи.

\subsection{Линейно разширение и топологично сортиране}

\begin{keythm}[Влагане в пълна наредба]
Нека $A$ е непразно крайно множество и $R$ е частична наредба над $A$. Съществува пълна наредба $S$ над $A$, за която
\[
R\subseteq S.
\]
\end{keythm}

\begin{proof}
Индукция по $|A|$. За едноелементно множество твърдението е ясно. Избираме минимален елемент $a$; той съществува по предходната теорема. По индукция ограничението на $R$ върху $A\setminus\{a\}$ има линейно разширение. Поставяме $a$ пред всички елементи на това разширение. Така получаваме пълна наредба. Няма отношение $bRa$ за $b\ne a$ по минималността, а всички останали отношения на $R$ се запазват по индукционната хипотеза. Следователно новата наредба съдържа $R$.
\end{proof}

\begin{defn}[Линейно разширение]
Релацията $S$ се нарича линейно разширение на $R$. Построяването му е известно като топологично сортиране.
\end{defn}

\begin{verbatim}
Вход: крайно множество A и частична наредба R над A
Изход: редица B, задаваща линейно разширение

i <- 1
докато A не е празно:
    избери минимален елемент a на A спрямо R
    B[i] <- a
    изтрий a от A и ограничи R върху оставащите елементи
    i <- i + 1
върни B
\end{verbatim}

Идеята за коректност е следната. Във всяка непразна стъпка оставащото крайно частично наредено множество има минимален елемент, следователно алгоритъмът може да продължи. Ако елемент $a$ предшества $b$ по $R$, $b$ не може да бъде избран преди $a$, защото тогава $b$ не би бил минимален сред оставащите елементи. Следователно редицата запазва всички отношения от $R$ и задава пълна наредба, съдържаща $R$.

\section{Функции}

\begin{keydefn}[Частична функция]
Нека $A$ и $B$ са множества. Релация $f\subseteq A\times B$ е частична функция от $A$ в $B$, ако за всеки $a\in A$ има най-много един $b\in B$, за който $(a,b)\in f$:
\[
(\forall a\in A)(\forall b_1,b_2\in B)\,
\bigl((a,b_1)\in f\land(a,b_2)\in f\Rightarrow b_1=b_2\bigr).
\]
\end{keydefn}

При частична функция не е задължително всеки елемент на $A$ да има стойност.

\begin{keydefn}[Тотална функция]
Частичната функция $f\subseteq A\times B$ е тотална функция от $A$ в $B$, ако
\[
\operatorname{dom}(f)=A.
\]
Еквивалентно:
\[
(\forall a\in A)(\exists!b\in B)\ (a,b)\in f.
\]
Пише се $f:A\to B$, а вместо $(a,b)\in f$ се пише $f(a)=b$.
\end{keydefn}

\begin{keydefn}[Инекция, сюрекция и биекция]
Нека $f:A\to B$ е тотална функция.

\begin{itemize}
\item $f$ е \emph{инекция}, ако
\[
(\forall a_1,a_2\in A)\,
(f(a_1)=f(a_2)\Rightarrow a_1=a_2).
\]

\item $f$ е \emph{сюрекция}, ако
\[
(\forall b\in B)(\exists a\in A)\ f(a)=b.
\]

\item $f$ е \emph{биекция}, ако е едновременно инекция и сюрекция.
\end{itemize}
\end{keydefn}

Инекцията не допуска различни аргументи с една и съща стойност. Сюрекцията изисква всяка стойност от $B$ да бъде достигната. Биекцията установява съответствие един към един между елементите на $A$ и $B$.

\section{Крайни и изброими множества. Принцип на Дирихле}

\begin{keydefn}[Крайно множество и кардиналност]
Множество $A$ е крайно, ако $A=\varnothing$ или съществува $n\in\mathbb N$, за което има биекция
\[
A\longrightarrow \{0,1,\ldots,n-1\}.
\]
Числото $n$ се нарича кардиналност на $A$ и се означава с $|A|=n$. За празното множество $|\varnothing|=0$.
\end{keydefn}

\begin{keydefn}[Изброимо безкрайно множество]
Множество $A$ е изброимо безкрайно, ако съществува биекция между $A$ и $\mathbb N$.
\end{keydefn}

\begin{keythm}[Принцип на Дирихле]
Нека $X$ и $Y$ са крайни множества и
\[
|X|>|Y|.
\]
Тогава за всяка тотална функция $f:X\to Y$ съществуват различни $a,b\in X$, за които
\[
f(a)=f(b).
\]
Следователно не съществува инекция от $X$ в $Y$.
\end{keythm}

\begin{proof}
Ако никои два различни аргумента нямат един и същ образ, всяко влакно $f^{-1}(\{y\})$ има най-много един елемент. Тези влакна са непресичащи се и покриват $X$, понеже $f$ е тотална. Следователно
\[|X|=\sum_{y\in Y}|f^{-1}(\{y\})|\le |Y|,\]
което противоречи на хипотезата. По същия начин, ако всяка от $n$ кутии съдържа най-много $k$ обекта, общият брой е най-много $kn$, което доказва и обобщената форма.
\end{proof}

Обобщената форма гласи: ако $kn+1$ обекта се разпределят в $n$ кутии, то поне в една кутия има повече от $k$ обекта.

\begin{example}
При разпределяне на $13$ студенти в $12$ месеца според месеца им на раждане поне двама студенти са родени в един и същ месец. Това е пряко приложение на принципа на Дирихле.
\end{example}

\section{Примерни задачи}

\begin{example}
Нека
\[
A=\{1,2,3\},\qquad B=\{2,3,4\}.
\]
Тогава
\[
A\cup B=\{1,2,3,4\},\qquad
A\cap B=\{2,3\},
\]
\[
A\setminus B=\{1\},\qquad
A\triangle B=\{1,4\}.
\]
\end{example}

\begin{example}
Нека $A=\{1,2\}$ и $B=\{a,b\}$. Тогава
\[
A\times B=\{(1,a),(1,b),(2,a),(2,b)\}.
\]
Релацията
\[
R=\{(1,a),(2,b)\}\subseteq A\times B
\]
е тотална функция от $A$ към $B$. Тя е биекция, защото различните елементи на $A$ имат различни образи и всеки елемент на $B$ е образ на елемент от $A$.
\end{example}

\begin{example}
Нека $A=\{1,2,3,6\}$ и над $A$ е зададена релацията делимост:
\[
aRb\quad\Longleftrightarrow\quad a\mid b.
\]
Тя е частична наредба. Елементите $2$ и $3$ са несравними, защото нито $2\mid 3$, нито $3\mid 2$. Минимален елемент е $1$, а максимален елемент е $6$.
\end{example}

\begin{example}
Нека $R$ над $\mathbb Z$ е зададена с
\[
aRb\quad\Longleftrightarrow\quad a-b\text{ е четно}.
\]
Релацията е рефлексивна, симетрична и транзитивна, следователно е релация на еквивалентност. Тя разделя $\mathbb Z$ на два класа: клас на четните и клас на нечетните цели числа.
\end{example}

\section{Изпитен фокус}

\begin{itemize}
\item Да се формулират аксиомата за обема, схемата за отделянето, аксиомата за степенното множество и идеята за индуктивно генерирано множество.
\item Да се обяснят слабата и пълната математическа индукция и да се проведе индукционно доказателство с база и индуктивна стъпка.
\item Да се дефинират обединение, сечение, разлика, симетрична разлика, допълнение и степенно множество; да се използват законите на Де Морган.
\item Да се дефинират наредена двойка, наредена $n$-орка, декартово и обобщено декартово произведение.
\item Да се дефинира $n$-местна релация, както и домейн и множество от стойности на бинарна релация.
\item Да се формулират рефлексивност, антирефлексивност, симетричност, антисиметричност, силна антисиметричност, асиметричност и транзитивност.
\item Да се дефинират релация на еквивалентност, клас на еквивалентност и разбиване на множество.
\item Да се разграничават частична и пълна наредба; да се определят минимални и максимални елементи и да се интерпретира диаграма на Хассе.
\item Да се възпроизведе идеята на топологичното сортиране: многократно избиране и премахване на минимален елемент, докато се получи линейно разширение.
\item Да се дефинират частична и тотална функция, инекция, сюрекция и биекция.
\item Да се дефинират крайно множество, кардиналност и изброимо безкрайно множество.
\item Да се формулира принципът на Дирихле и да се разпознава приложението му като доказателство, че две различни стойности имат еднакъв образ.
\end{itemize}
